\documentclass[11pt]{amsart}
\usepackage{amssymb}

\newtheorem{theorem}{Theorem}[section]
\newtheorem{lemma}[theorem]{Lemma}
\newtheorem{proposition}[theorem]{Proposition}
\newtheorem{corollary}[theorem]{Corollary}
\theoremstyle{definition}
\newtheorem{definition}[theorem]{Definition}

\newcommand{\Z}{\mathbb{Z}}
\newcommand{\Q}{\mathbb{Q}}

\title{Divisibility and rational roots of integer polynomials}
\author{Marta K.~Olsberg}
\address{Department of Mathematics, Lindqvist College, 41 Harbour Road, Kalmarsund}
\email{m.olsberg@lindqvist.example.edu}
\date{}

\begin{document}

\begin{abstract}
We give a short, self-contained account of the elementary divisibility
properties of polynomials with integer coefficients. The main result is the
classical fact that every rational root of a monic integer polynomial is an
integer dividing the constant term. Along the way we record several
consequences of the basic congruence $a-b \mid P(a)-P(b)$, including a
criterion for the absence of integer roots and a remark on periodic points.
\end{abstract}

\maketitle

\section{Introduction}

Polynomials with integer coefficients enjoy a rigid arithmetic structure:
their values at integers are constrained by divisibility relations that do
not hold for arbitrary real polynomials. The most basic of these is the
observation that $a-b$ divides $P(a)-P(b)$ whenever $a,b\in\Z$ and
$P\in\Z[x]$. Combined with unique factorisation in $\Z$, it yields the
\emph{rational root theorem}, which in turn shows that the rational roots of
a monic integer polynomial are integers.

These facts are standard and appear in many textbooks, but they are usually
scattered over several chapters. The purpose of this note is to collect them
in one place, with complete proofs, in a form convenient for teaching and for
quick reference. Our main result is Theorem~\ref{thm:main}. In
Section~\ref{sec:app} we combine it with Eisenstein's criterion and derive two
simple non-existence results for rational roots.

\section{Definitions and main results}

\begin{definition}\label{def:basic}
Throughout, $P\in\Z[x]$ denotes a nonzero polynomial, written
$P(x)=a_nx^n+a_{n-1}x^{n-1}+\dots+a_1x+a_0$ with $a_i\in\Z$ and $a_n\neq 0$,
so that $n=\deg P$. We call $a_n$ the \emph{leading coefficient} and
$a_0=P(0)$ the \emph{constant term}; $P$ is \emph{monic} if $a_n=1$. A
\emph{root} of $P$ is a number $r$ with $P(r)=0$. A rational number $r$ is
\emph{written in lowest terms} as $r=p/q$ if $p,q\in\Z$, $q\geq 1$ and
$\gcd(p,q)=1$; every $r\in\Q$ has exactly one such representation.
\end{definition}

\begin{lemma}\label{lem:diff}
Let $P\in\Z[x]$ and $a,b\in\Z$. Then $a-b$ divides $P(a)-P(b)$.
\end{lemma}

\begin{proof}
For every $k\geq 1$ we have the factorisation
\[
a^k-b^k=(a-b)\sum_{j=0}^{k-1}a^{j}b^{k-j},
\]
so $a-b \mid a^k-b^k$. Since
$P(a)-P(b)=\sum_{k=1}^{n}a_k\,(a^k-b^k)$, the claim follows by linearity.
\end{proof}

\begin{corollary}\label{cor:cong}
Let $m\geq 1$ and let $P\in\Z[x]$ be monic. If $a,b\in\Z$ satisfy
$a\equiv b \pmod m$, then $P(a)\equiv P(b)\pmod m$.
\end{corollary}

\begin{proof}
By hypothesis $m\mid a-b$, and by Lemma~\ref{lem:diff} $a-b\mid P(a)-P(b)$.
Hence $m\mid P(a)-P(b)$ by transitivity of divisibility.
\end{proof}

\begin{lemma}\label{lem:noroot}
Let $m\geq 2$ and $P\in\Z[x]$. If $m\nmid P(k)$ for every
$k\in\{1,2,\dots,m-1\}$, then $P$ has no integer root.
\end{lemma}

\begin{proof}
Suppose $r\in\Z$ satisfies $P(r)=0$, and let $k\in\{1,\dots,m-1\}$ be the
residue of $r$ modulo $m$. Then $m\mid r-k$, and by Lemma~\ref{lem:diff},
$r-k \mid P(r)-P(k)=-P(k)$. Thus $m\mid P(k)$, contrary to the hypothesis.
\end{proof}

\begin{lemma}\label{lem:unit}
Let $P\in\Z[x]$ and $a,b\in\Z$. If $|P(a)-P(b)|=1$, then $|a-b|=1$.
\end{lemma}

\begin{proof}
By Lemma~\ref{lem:diff}, $a-b$ divides $P(a)-P(b)=\pm 1$, so
$|a-b|\leq 1$. If $a=b$ then $P(a)-P(b)=0$, which is impossible; hence
$|a-b|=1$.
\end{proof}

\begin{lemma}\label{lem:introot}
Let $P\in\Z[x]$. Every integer root $r$ of $P$ divides the constant term
$a_0$; in particular $|r|\leq |a_0|$.
\end{lemma}

\begin{proof}
Apply Lemma~\ref{lem:diff} with $a=r$ and $b=0$: $r\mid P(r)-P(0)=-a_0$.
A nonzero divisor of an integer is at most that integer in absolute value.
\end{proof}

\begin{proposition}\label{prop:cycle}
Let $P\in\Z[x]$ and let $a,b,c$ be pairwise distinct integers with
$P(a)=b$, $P(b)=c$ and $P(c)=a$. Then $|a-b|=|b-c|=|c-a|$.
\end{proposition}

\begin{proof}
By Lemma~\ref{lem:diff}, $a-b\mid P(a)-P(b)=b-c$, $b-c\mid P(b)-P(c)=c-a$
and $c-a\mid P(c)-P(a)=a-b$. All three differences are nonzero, so
$|a-b|\leq|b-c|\leq|c-a|\leq|a-b|$, and equality holds throughout.
\end{proof}

\begin{theorem}[Rational root theorem]\label{thm:rrt}
Let $P\in\Z[x]$ have degree $n\geq 1$, and let $r=p/q\in\Q$, written in
lowest terms, be a root of $P$. Then $p\mid a_0$ and $q\mid a_n$.
\end{theorem}

\begin{proof}
Multiplying $P(p/q)=0$ by $q^n$ gives
$\sum_{i=0}^{n}a_i\,p^{i}q^{n-i}=0$. Every term with $i\geq 1$ is divisible
by $p$, hence $p\mid a_0q^n$; since $\gcd(p,q^n)=1$ we get $p\mid a_0$.
Similarly every term with $i\leq n-1$ is divisible by $q$, hence
$q\mid a_np^n$, and $\gcd(q,p^n)=1$ gives $q\mid a_n$.
\end{proof}

\begin{theorem}\label{thm:main}
Let $P\in\Z[x]$ be monic of degree $n\geq 1$, and let $r\in\Q$ be a root of
$P$. Then $r\in\Z$ and $r$ divides the constant term $P(0)$.
\end{theorem}

\begin{proof}
Write $r=p/q$ in lowest terms. By Theorem~\ref{thm:rrt}, $q\mid a_n=1$, and
since $q\geq 1$ this forces $q=1$, i.e.\ $r=p\in\Z$. Theorem~\ref{thm:rrt}
also gives $p\mid a_0=P(0)$.
\end{proof}

\section{Applications}\label{sec:app}

We recall a classical irreducibility criterion.

\begin{theorem}[Eisenstein's criterion {\cite[Ch.~IV, Thm.~3.1]{Lang}}]
\label{thm:eis}
Let $P=a_nx^n+\dots+a_0\in\Z[x]$ with $n\geq 1$, and let $p$ be a prime
such that $p\nmid a_n$, $p\mid a_i$ for all $0\leq i<n$, and
$p^2\nmid a_0$. Then $P$ is irreducible in $\Z[x]$.
\end{theorem}

For instance, for every prime $p$ and every $n\geq 1$ the polynomial
$x^n-p$ is irreducible. The absence of rational roots, which is all we need
below, can also be read off directly from Theorem~\ref{thm:main}.

\begin{corollary}\label{cor:xnp}
Let $p$ be a prime and $n\geq 2$. Then $x^n-p$ has no rational root.
\end{corollary}

\begin{proof}
By Theorem~\ref{thm:main}, a rational root $r$ would be an integer dividing
$p$, so $r\in\{\pm1,\pm p\}$. But then $|r|^n\in\{1,p^n\}$, and neither
equals $p$ because $n\geq 2$.
\end{proof}

\begin{corollary}\label{cor:parity}
Let $P\in\Z[x]$ be monic of degree $n\geq 1$ with $P(0)$ and $P(1)$ both
odd. Then $P$ has no rational root.
\end{corollary}

\begin{proof}
By Theorem~\ref{thm:main}, any rational root of $P$ is an integer, and by
Lemma~\ref{lem:noroot} with $m=2$, $P$ has no integer root.
\end{proof}

\begin{thebibliography}{9}
\bibitem{Lang} S.~Lang, \emph{Algebra}, revised third edition, Graduate
Texts in Mathematics 211, Springer, New York, 2002.
\end{thebibliography}

\end{document}
